\documentclass[11pt,letterpaper]{article}
\usepackage[margin=1in]{geometry}
\usepackage{booktabs}
\usepackage{longtable}
\usepackage{hyperref}
\usepackage{xcolor}
\usepackage{enumitem}
\usepackage{parskip}

\title{Independent Replication Report --- BVBRC-78\\
\large \textit{Enterococcus faecium} Bacteriophage vB\_EfaH\_163 (Pradal et al., \textit{Viruses} 2023)}
\author{Independent computational replication\\Ollie / OpenClaw REPLICATE-PROJECT wave BVBRC}
\date{Verdict: \textbf{PARTIAL}}

\begin{document}
\maketitle

\section{Paper Under Replication}

\begin{description}[leftmargin=1.3cm,style=nextline]
  \item[Title] \textit{Enterococcus faecium} Bacteriophage vB\_EfaH\_163, a New Member of the Herelleviridae Family, Reduces the Mortality Associated with an \textit{E. faecium} vanR Clinical Isolate in a \textit{Galleria mellonella} Animal Model
  \item[Authors] Pradal I, Casado A, del Rio B, Rodr\'iguez-Lucas C, Fern\'andez M
  \item[Venue] \textit{Viruses} 15(1): 179 (2023) \textbullet{} PMID 36680219 \textbullet{} DOI 10.3390/v15010179 \textbullet{} PMC9860891
  \item[Deposited data] ENA WGS accession \texttt{CAJDKA010000002.1} (phage genome contig only). The host isolate genome and the raw sequencing reads are \emph{not} deposited.
\end{description}

\section{Claims Extracted From The Paper}

Thirteen discrete claims were extracted (C-1--C-13): eight genomic/comparative (computationally testable in principle), one packaging claim requiring raw reads (C-9), three wet-lab claims out of scope for computational replication (C-10--C-12), and one host-genome claim blocked by the authors not depositing the host isolate (C-13).

\begin{center}
\small
\begin{longtable}{clcccp{4.8cm}l}
\toprule
\# & Claim (abridged) & Type & Testable? & Tested? & Independent result & Verdict\\
\midrule
\endhead
C-1  & Genome $=150{,}836$ bp                     & genomic     & \checkmark & \checkmark & 150,836 bp exact & AGREE\\
C-2  & GC $\approx 37\%$                           & genomic     & \checkmark & \checkmark & 37.04\% & AGREE\\
C-3  & 186 ORFs                                    & genomic     & \checkmark & \checkmark & Prodigal (meta): 183 ORFs & AGREE ($\pm 2\%$)\\
C-4  & 21 tRNAs                                    & genomic     & \checkmark & \checkmark & ARAGORN: 21 tRNAs & AGREE (exact)\\
C-5  & No AMR/virulence genes                      & safety      & \checkmark & \checkmark & Abricate 7 DBs: 0 hits & AGREE\\
C-6  & Lytic (no lysogeny genes)                   & genomic     & \checkmark & \checkmark & BLASTp vs curated integrase/cI: 0 hits & AGREE\\
C-7  & Herelleviridae / Brockvirinae / Schiekvirus & taxonomy    & \checkmark & \checkmark & Whole-genome BLASTn + MCP UPGMA agree & AGREE\\
C-8  & Top hits iF6/EfV12-phi1/EFDG1 $\sim 98\%$   & comparative & \checkmark & \checkmark & Same three phages top; 93.8--96.5\% & AGREE (dir.)\\
C-9  & Long DTRs (PhageTerm)                       & genomic     & \checkmark & $\times$   & Raw reads not deposited & NOT TESTED\\
C-10 & Host range 51\%, 16 VRE                     & wet lab     & $\times$   & --         & Out of scope & N/T\\
C-11 & Burst size $\approx$155 PFU, latent 60 min  & wet lab     & $\times$   & --         & Out of scope & N/T\\
C-12 & \textit{Galleria} mortality reduction       & wet lab     & $\times$   & --         & Out of scope & N/T\\
C-13 & VR-13 carries van cluster                   & genomic     & \checkmark & $\times$   & Host genome not deposited & NOT TESTED\\
\bottomrule
\end{longtable}
\end{center}

\textbf{Score:} 7 tested / 8 accessible-in-principle; 7/7 tested claims agree with the paper (100\%). Two accessible-in-principle genomic claims (C-9, C-13) are blocked by undeposited data.

\section{Methods (condensed)}

All analyses were performed on macOS using Homebrew-installed BLAST+ 2.16, Prodigal 2.6.3 (via \texttt{pyrodigal} 3.7.1), ARAGORN 1.2.41, Abricate 1.0.1, and MAFFT 7 (the last one segfaulted, see below). Python was 3.14 with BioPython 1.87. The phage FASTA was downloaded directly from ENA (\texttt{CAJDKA010000002.1}); five Herelleviridae comparators (iF6 \texttt{NC\_029009}, EfV12-phi1 \texttt{MH880817}, EFDG1 \texttt{NC\_047796.1}, EFP01 \texttt{MT909815.1}, MDA2 \texttt{MW633168.1}) and two Siphoviridae outgroups (\texttt{NC\_031260}, \texttt{MK360024}) were fetched via NCBI eFetch. Whole-genome BLASTn was computed pairwise with weighted-average \%identity across all HSPs. A curated 7-protein lysogeny-marker set (integrases from $\lambda$, $\phi 80$, P22, $\phi$Sa3int, L54a plus $\lambda$ cI and P22 cI) was used to screen the vB\_EfaH\_163 proteome by BLASTp at $E < 10^{-5}$. Major-capsid-protein phylogeny was built by extracting putative MCPs via BLASTp against EFDG1 major head protein \texttt{YP\_009218324.2}, followed by all-vs-all pairwise identities (BioPython \texttt{PairwiseAligner}) and UPGMA on $1 - \text{pid}$ distances. Four LLM judges (Argo proxy: \texttt{gpt-5.2}, \texttt{gemini-2.5-pro}, \texttt{claude-sonnet-4.6}, \texttt{gpt-5.4}) scored the evidence packet at temperature 0.1 with an enforced verdict vocabulary; two additional Argo opus endpoints (4.7, 4.8) returned HTTP 502 during the run.

\section{Results vs Paper (numeric)}

\begin{center}
\small
\begin{tabular}{lccc}
\toprule
Metric & Paper & This work & $\Delta$\\
\midrule
Genome length (bp)        & 150{,}836 & 150{,}836 & 0 \\
GC \%                     & $\sim$37   & 37.04    & $\sim$0 \\
ORF count                 & 186       & 183       & $-3$ (caller variation) \\
tRNA count                & 21        & 21        & 0 \\
AMR gene hits             & 0         & 0 (7 DBs) & 0 \\
Virulence gene hits       & 0         & 0 (2 DBs) & 0 \\
Lysogeny markers          & none      & 0 ($E<10^{-5}$) & 0 \\
BLASTn vs iF6             & $\sim$98\% & 96.5\%    & $-1.5$ pp \\
BLASTn vs EFDG1           & $\sim$98\% & 93.8\%    & $-4.2$ pp \\
MCP tree topology (Fig 4) & recovered & recovered (iF6 100\%; MDA2 outlier) & agree \\
\bottomrule
\end{tabular}
\end{center}

\section{LLM Judge Panel (Argo :44497)}

\begin{center}
\begin{tabular}{lccc}
\toprule
Judge & Verdict & Coverage & Agreement \\
\midrule
\texttt{argo:gpt-5.2}            & PARTIAL    & 8/9  & 8/8 \\
\texttt{argo:gemini-2.5-pro}     & REPLICATED & 8/8  & 8/8 \\
\texttt{argo:claude-sonnet-4.6}  & PARTIAL    & 7/10 & 7/7 \\
\texttt{argo:gpt-5.4}            & PARTIAL    & 7/10 & 7/7 \\
\bottomrule
\end{tabular}
\end{center}

Majority verdict of the four responsive judges: \textbf{PARTIAL} (3 of 4). Gemini alone rated the paper REPLICATED because it computed coverage over accessible-in-principle claims only; the other three penalized the two blocked-by-missing-data claims.

\section{Genuine Critique}

This is the honest, self-critical section --- what actually worked, what did not, and where the replication is weaker than a single-number verdict suggests.

\subsection*{What replicated cleanly}
The paper's \emph{static} genomic claims all reproduce with essentially no ambiguity. Genome length is exact to the base pair, tRNA count is exact, GC agrees to two decimal places, and safety-relevant screens (AMR, virulence, lysogeny) are unambiguously negative across a broad panel of databases and curated reference proteins. Taxonomic placement (Herelleviridae / Brockvirinae / Schiekvirus, closest to iF6/EFP01/EfV12-phi1/EFDG1) recovers not only the paper's assignment but also its Fig 4 topology (MDA2 falls on the Kochikohdavirus branch as the paper reported). None of this required any tuning to reproduce.

\subsection*{What only replicated in the correct \emph{direction}}
The BLASTn \%identity numbers drifted downward by 1.5--4.2 percentage points against the paper's ``$\sim$98\%'' summary. This is almost certainly a tool-choice artefact: the paper used a whole-genome pipeline (VIRIDIC/pyani or similar) that averages differently than megablast's weighted-HSP average, and rounding hides several percentage points at the reporting resolution. The \emph{ordering} of top hits is unchanged. We flag this because ``96.5\% vs 98\%'' looks reassuring in a table but a naive reader might interpret a systematic $-3$ pp drift as evidence of a data-quality issue when it is really a metric-definition issue.

\subsection*{What did \emph{not} replicate --- and why the verdict is PARTIAL not REPLICATED}
\begin{itemize}[leftmargin=1.2cm]
  \item \textbf{C-9 (long direct terminal repeats via PhageTerm).} PhageTerm is a read-mapping tool; it needs raw reads to infer packaging strategy from read-end pile-ups. The authors deposited only the assembled contig, so this claim is \emph{computationally unfalsifiable from public data}. A reader who wants to independently confirm the packaging strategy cannot.
  \item \textbf{C-13 (VR-13 carries a van cluster).} The paper's therapeutic story hinges on the host being a bona-fide vancomycin-resistant isolate carrying a \textit{van} gene cluster. Because the host isolate genome was not deposited, we cannot BLAST a \textit{vanA}/\textit{vanB} reference against it. The phenotype (vancomycin MIC in the paper) is reported but the genotype is not independently confirmable.
  \item \textbf{C-10, C-11, C-12 (wet-lab).} Out-of-scope for computational replication in the strict sense, but worth flagging: the paper's clinical translation narrative (host range 51\%, one-step growth, and above all the \textit{Galleria mellonella} mortality reduction) is entirely wet-lab, entirely non-independently-verified here, and depends on strain-specific reagents that are not distributable.
\end{itemize}

\subsection*{Method weaknesses in \emph{our} replication}
\begin{enumerate}[leftmargin=1.4cm]
  \item \textbf{Prodigal in meta mode is not a curated ORF call.} We report 183 vs the paper's 186. The paper used RAST + PATRIC + manual BLAST curation, which is the correct pipeline for a phage annotation and will find short/frameshifted/overlapping ORFs Prodigal misses. Our ``$\pm 2\%$'' verdict is charitable; a stricter reviewer might call this UNVERIFIED-BY-INDEPENDENT-METHOD rather than AGREE.
  \item \textbf{MAFFT segfaulted on Homebrew macOS build.} We fell back to BioPython \texttt{PairwiseAligner} + UPGMA, which is defensible for a 6-protein MCP tree but is not the same as a proper MSA + ML tree. The topology recovered matches the paper's Fig 4, so the qualitative conclusion holds, but branch lengths in our tree should not be over-interpreted.
  \item \textbf{Lysogeny screen used 7 curated markers, not an HMM sweep.} We chose $\lambda$/$\phi 80$/P22/Sa3int/L54a integrases plus $\lambda$ and P22 cI. A more thorough screen would use PHASTER/PHASTEST HMMs or the full VOG lysogeny catalogue. A true integrase from an unusual family could in principle escape our screen; the paper's own screening pipeline is similarly narrow so we consider this a shared limitation, not a contradiction.
  \item \textbf{LLM-judge panel had 2/6 endpoints (Argo opus 4.7 and 4.8) return HTTP 502.} We substituted \texttt{claude-sonnet-4.6} and \texttt{gpt-5.4} to keep the panel at 4 responsive judges. This is defensible diversity but is not the pre-registered panel.
  \item \textbf{No sensitivity analysis on BLASTn identity metric.} We used one metric (weighted-average HSP \%identity). A sensitivity analysis using VIRIDIC (paper's likely tool) or pyani ANIm would let us convert ``directional agreement'' into a proper numeric agreement.
\end{enumerate}

\subsection*{What we do \emph{not} claim}
We do not claim to have replicated the paper's therapeutic conclusion. Our verdict PARTIAL applies specifically to the \emph{computational/genomic} substrate underneath that conclusion. The clinical relevance --- that vB\_EfaH\_163 usefully reduces mortality of a VRE isolate in an insect model, let alone in a mammalian model or a patient --- is entirely outside what we tested.

\section{Verdict}

\textbf{PARTIAL.} Every accessible computational claim reproduced (7/7 tested, 100\% agreement with only minor numerical drift on BLASTn \%id). The verdict is not REPLICATED because two claims that are computational in nature (C-9 packaging, C-13 host van cluster) are blocked by the authors not depositing raw reads or the host isolate genome. This is a data-sharing gap, not a factual disagreement. The paper's genomic substrate is trustworthy; the therapeutic conclusion is out of scope for this replication.

\section{Files}

\begin{itemize}[noitemsep]
  \item \texttt{report/REPORT.md}, \texttt{report/REPORT.tex} --- this report
  \item \texttt{report/brief.md} --- one-paragraph summary
  \item \texttt{report/attempt\_log.md} --- chronological log
  \item \texttt{report/artifact\_harvest.md} --- every public artefact fetched
  \item \texttt{report/evidence/} --- Abricate outputs, BLASTn summary, MCP tree, ARAGORN, Prodigal calls, judge input + verdicts
  \item \texttt{work/} --- raw downloads, intermediates, analysis scripts
\end{itemize}

\end{document}
